\documentclass{article}

\usepackage{booktabs}
\usepackage{graphicx}
\usepackage{hyperref}

\title{Calibrated Anchor Replay: Reducing Recency Bias and Representation Drift in Class-Incremental Learning}
\author{Utkarsh Rajput}
\date{}

\begin{document}
\maketitle

\begin{abstract}
Class-incremental continual learning requires models to acquire new classes
without overwriting earlier decision boundaries. We study whether replay methods
can be strengthened with lightweight logit anchors, feature anchors, class
prototypes, and post-task calibration while preserving fixed memory budgets. The
central claim is evaluated through memory-accuracy Pareto curves under matched
Split CIFAR-10, Split CIFAR-100, and TinyImageNet protocols.
\end{abstract}

\section{Introduction}
State the plasticity-stability problem, why replay remains a strong baseline,
and why matched memory budgets are necessary for fair claims.

\section{Method}
Define Calibrated Anchor Replay. The method stores exemplars, labels, teacher
logits, penultimate features, and prototype statistics. Training combines current
cross-entropy, replay cross-entropy, logit matching, feature anchoring, and
prototype anchoring. After each task, a temperature and bias calibration head is
fit over balanced memory.

\section{Experimental Protocol}
Report full-data Split CIFAR-10, Split CIFAR-100, and TinyImageNet. Compare
baseline, joint cumulative-data training, replay, DER++, ER-ACE, A-GEM, GDumb,
iCaRL-style replay, BiC-style calibration, X-DER-lite, and CAR under matched
memory budgets.

\section{Results}
Insert generated tables and figures from \texttt{docs/paper/assets}. Claims must
be restricted to matched protocols.

\section{Ablations}
Remove balanced memory, logit anchors, feature anchors, prototype anchors, and
calibration one at a time.

\section{Limitations}
Discuss compute budget, dataset scope, external protocol differences, and any
cases where CAR fails to beat strong baselines.

\end{document}
